\documentclass[11pt]{article}
\usepackage{icmlko}
\begin{document}

\title{긁기 전에 상금을 쟀더니 내가 적은 규칙이 틀렸다 --- 상한을 추정값으로 쓴 사전 등록을, 발동했는데도 안 따랐다}
\author{월드모델 실험실}
\date{노트 548 · 2026년 8월 4일}
\maketitle

\begin{abstract}
대기 중이던 갈래 ``모바일 $\cdot$ 게임 검색 덮음 $30\!\sim\!41\%$ 제목 언어
한계''를 열었다. 펀딩은 덮음이 $1.0\%$ 인데 수집기에 \textbf{남은 요청이
$588$ 개} 있다. \textbf{긁기 전에 상금부터 쟀다} --- 검색 축을 통째로 끄면
얼마나 잃나.

\begin{center}
\begin{tabular}{lrrr}
\hline
도메인 & 유보 & 검색 덮음 & \textbf{끄면} \\
\hline
도서 & $163$ & $39.9\%$ & $\mathbf{-0.1654}$ \\
게임 & $180$ & $63.3\%$ & $-0.0515$ \\
\textbf{애니} & $606$ & $36.0\%$ & $\mathbf{-0.0003}$ \\
\textbf{펀딩} & $529$ & $\mathbf{1.0\%}$ & $-0.0016$ \\
\textbf{만화} & $258$ & $0.8\%$ & $\mathbf{+0.0099}$ \\
\textbf{팝업} & $65$ & $\mathbf{78.7\%}$ & $\mathbf{+0.0374}$ ($12/12$) \\
\hline
\textbf{판} & & & $\mathbf{-0.0204}$ \\
\hline
\end{tabular}
\end{center}

\textbf{덮음과 값 사이에 관계가 없다} --- 순위상관 $+0.191$($n=11$)이고
덮음 $30\%$ 이상 여섯만 보면 $\mathbf{-0.143}$ 이다. 덮음 $79\%$ 인 둘이
각각 해롭거나 $0$ 이다.
\end{abstract}

\begin{figure}[t]
\centering
\includegraphics[width=\textwidth]{figs/f1.pdf}
\caption{\textbf{가} --- 스무 자 절벽. \textbf{나} --- 결측이 정보다.
\textbf{다} --- 덮음과 값이 안 맞는다.}
\end{figure}

\section{갈래의 전제가 틀렸다}

``제목 언어 한계''가 아니다. \textbf{데이터랩은 영문 질의에 값을 준다}
(노트 $237$ 이 이미 확정했다) --- 게임에서 값을 받은 $347$ 건 중
\textbf{$56\%$ 가 영문 제목}이고, 원천에 아예 없는 $135$ 건은 한글 제목이
$5.9\%$ 다. \textbf{언어가 막은 것이 아니라 안 긁은 것}이다.

\begin{center}
\begin{tabular}{lrrr}
\hline
질의어 & 값 & 결과 없음 & \textbf{빈손 비율} \\
\hline
$<20$ 자 & $4{,}152$ & $3{,}656$ & $46.8\%$ \\
\textbf{$=20$ 자(잘림)} & $53$ & $2{,}133$ & $\mathbf{97.6\%}$ \\
\hline
\end{tabular}
\end{center}

\textbf{전체 빈손의 $36.8\%$ 가 스무 자 잘림}이다. 그리고 이것은
\textbf{이미 아는 결함이다} --- 노트 $490 \cdot 491$ 이 회귀 불연속으로
잡았고(정확히 $20$ 자 $18.3\%$ 대 $21$ 자 $0.6\%$, 승산비 $29.2$)
``\_cap'' 을 낱말 경계로 고쳤다. \textbf{그런데 다시 안 긁어서} 저장 파일에
옛 $20$ 자 질의어가 $2{,}179$ 건 남아 있다.

\begin{quote}
\textbf{null 은 수집 실패가 아니다.} 같은 질의어가 한 파일에선 값이 있고
다른 파일에선 없는 경우는 \textbf{$49$ 개 $\cdot$ 전체 null 의 $0.9\%$}
뿐이다. 재시도로 되찾을 것이 아니다.
\end{quote}

\section{결측 자체가 정보다}

\textbf{전 도메인에서 ``검색에 안 잡힘''이 라벨과 양수}다(연도 통제) ---
애니 $+0.3522$ $\cdot$ 도서 $+0.1808$ $\cdot$ 웹툰 $+0.1600$ $\cdot$ 게임
$+0.1590$ $\cdot$ 모바일 $+0.1233$ $\cdot$ 만화 $+0.1049$ $\cdot$ 펀딩
$+0.0978$ $\cdot$ 세계애니 $+0.0506$.

\begin{quote}
``\emph{검색량이 문턱 아래다}''가 곧 ``\emph{무명하다}''이고,
\texttt{zero\_is\_data} 표시자가 이미 그것을 갖고 있다. \textbf{덮음을
채우는 값은 보이는 것보다 작다} --- 이진 표시자를 눈금으로 바꾸는 것뿐이다.
\end{quote}

\section{내가 적은 규칙이 틀렸다}

사전 등록에 이렇게 적었다 --- ``판 손실 $\ge 0.010$ 이면 펀딩 수집을
한다''. 손실은 $\mathbf{-0.0204}$ 다. \textbf{규칙이 발동했다.}

\begin{quote}
\textbf{그런데 안 따른다.} 같은 문장에 ``이 값이 \emph{모든 검색 수집의
상한}이다''라고 적어 놓고, \textbf{그 상한을 추정값처럼 쓰는 결정 규칙}을
달았다. 상한이 문턱을 넘는 것은 ``넘을 수도 있다''이지 ``넘는다''가 아니다.
\end{quote}

그리고 갈라 보니 추론 자체가 안 선다.

\begin{itemize}
\item \textbf{덮음이 값을 예측 못 한다} --- 순위상관 $+0.191$, 덮음
$30\%$ 이상만 보면 $-0.143$. 덮음 $78.7\%$ 인 팝업은 \emph{해롭고}
$79.0\%$ 인 아이돌은 $0$ 이다.
\item \textbf{펀딩 자신의 값이 $-0.0016$}($3/12$, $0$ 을 문다)이다.
\item 가장 후한 가정 --- 덮음을 $14\%$ 까지 올리고(잘림 고친 뒤 히트율)
\emph{도서의 값을 그대로 얻는다}고 쳐도 --- 판 $\mathbf{+0.0083}$ 으로
문턱 아래다. 그 가정에 근거도 없다.
\end{itemize}

\textbf{펀딩 수집을 안 한다.} 규칙은 결과를 보고 고치는 것이 아니라
\emph{잘못 쓴 양}을 바로잡는 것이다 --- 상한과 추정값은 다른 것이고,
그 구분은 사전 등록 문장 안에 이미 적혀 있었다.

\section{대신 실험이 후보를 냈다}

\begin{center}
\begin{tabular}{lrrr}
\hline
검색이 해로운 곳 & 덮음 & 끄면 & 무게 \\
\hline
\textbf{팝업} & $78.7\%$ & $\mathbf{+0.0374}$ ($12/12$) & $1.9\%$ \\
\textbf{만화} & $0.8\%$ & $\mathbf{+0.0099}$ ($10/12$) & $7.5\%$ \\
\textbf{KR 만화}(집 밖) & --- & $\mathbf{+0.0088}$ ($12/12$) & 자 \\
\hline
\end{tabular}
\end{center}

\textbf{애니는 이미 \texttt{DOMDROP} 에 들어 있다} --- 그래서 $-0.0003$
이다. 같은 처방이 두 곳 더 있을 수 있고, \textbf{집 밖 짝까지 같은 방향}
이라는 것이 강한 신호다. 노트 $551$ 에서 네 팔로 가른다.

\section*{조항}

\begin{quote}
\textbf{상한을 재는 실험에는 상한용 결정 규칙을 단다.} ``이 값이 상한이다''
와 ``이 값이 문턱을 넘으면 한다''를 같은 문단에 쓰면 안 된다.
\textbf{그리고 덮음은 값이 아니다} --- 자료를 더 모으는 결정은 덮음이 아니라
\emph{그 도메인에서 그 축이 버는 양}으로 한다.
\end{quote}

\section*{남은 것}

\begin{itemize}
\item \textbf{도서가 $-0.1654$} --- 검색이 그 도메인 $\rho$ 의 $41\%$ 를
낸다($0.4038 \to 0.2384$). 유보 $163$ 행뿐이라 무게는 $4.8\%$ 지만,
\emph{왜 도서만 그런지}는 안 봤다.
\item \textbf{스무 자 잘림 재수집} --- 대상 $2{,}015$ 건인데 노트 $491$ 의
추정대로면 되찾을 것이 $\sim\!170$ 행이고, 그 대부분이 만화($842$ 건
재질의 $\to$ 덮음 $0.8\% \to 2.4\%$)라 값이 없다. \textbf{안 한다.}
\item \textbf{시장팝업 덮음이 $0.0\%$} 인데 끄면 $-0.0103$ 이다 --- 제
검색 축이 없는데도 잃는다. 남의 학습이 바뀌어서 생기는 간접 효과다.
\end{itemize}

\end{document}
